%TC:macro \footnote [text]
%TC:macro \endnote [text]
%TC:incbib

%%%%%%%%%%%%%%%%%%%%%%%%%%%%%%%%%%%%%%%%%%%%%%%%%%%%%%%%%%%%%%%%%%%%%%%%%%%%%%%%
% Template for TISMIR Papers
% 2017 version, based on previous ISMIR conference template
%%%%%%%%%%%%%%%%%%%%%%%%%%%%%%%%%%%%%%%%%%%%%%%%%%%%%%%%%%%%%%%%%%%%%%%%%%%%%%%%

\documentclass{article}

%%%%%%%%%%%%%%%%%%%%%%%%%%%%%%%%%%%%%%%%%%%%%%%%%%%%%%%%%%%%%%%%%%%%%%%%%%%%%%%%
% Sample Document LaTeX packages
%%%%%%%%%%%%%%%%%%%%%%%%%%%%%%%%%%%%%%%%%%%%%%%%%%%%%%%%%%%%%%%%%%%%%%%%%%%%%%%%

\usepackage[utf8]{inputenc}
% To remove the line numbers for the camera-ready version of your paper, remove
% the `review` option from the tismir package
\usepackage[review]{tismir}
\usepackage{amsmath}
\usepackage{url}
\usepackage{graphicx}
\usepackage{booktabs}
\usepackage{lipsum}
\usepackage{placeins} 
\usepackage{afterpage} 
\usepackage{endnotes}
% ---- Highlighting packages (requires LuaLaTeX) ----
\usepackage[dvipsnames]{xcolor}
\usepackage{luacolor}
\usepackage{lua-ul}
\usepackage{xcolor}
\usepackage{float}

\definecolor{r2}{RGB}{0, 80, 200}  
\definecolor{r3}{RGB}{0, 90, 255}  
\definecolor{r4}{RGB}{0, 160, 80}   % round-3 (additional reviewer) changes
\newcommand{\rev}[1]{{\color{r2}#1}} % new-version (round-3) edits

% ---------------------------------------------------
\usepackage{xurl} % better URL breaking
\usepackage[hidelinks]{hyperref}

\def\UrlBreaks{\do\/\do-\do_\do.\do=\do?\do\&}

%%%%%%%%%%%%%%%%%%%%%%%%%%%%%%%%%%%%%%%%%%%%%%%%%%%%%%%%%%%%%%%%%%%%%%%%%%%%%%%%
% Title and Author information
%%%%%%%%%%%%%%%%%%%%%%%%%%%%%%%%%%%%%%%%%%%%%%%%%%%%%%%%%%%%%%%%%%%%%%%%%%%%%%%%

\title{Data-Driven Analysis of Musical Form and Harmonic Structure in AI-Generated Popular Music: A Case Study with Suno and Udio}
% Corpus-Scale Analysis of Form and Harmony in Human and AI-Generated Popular Music

%
\author{%
First Author1\thanks{University One, Campus Rd. 3, 12345 Mirland},%
~Second Author2\thanks{Research Institute Two, Campus Rd. 4, 12345 Flatland},%
~and Third Author3\protect\footnotemark[1]}

\date{}

%%%%%%%%%%%%%%%%%%%%%%%%%%%%%%%%%%%%%%%%%%%%%%%%%%%%%%%%%%%%%%%%%%%%%%%%%%%%%%%%
% Additional Paper Information
%%%%%%%%%%%%%%%%%%%%%%%%%%%%%%%%%%%%%%%%%%%%%%%%%%%%%%%%%%%%%%%%%%%%%%%%%%%%%%%%

% Article Type - Uncomment and modify, if necessary.
% Accepted values: research, overview, and dataset
% \type{research}

% Citation in First Page
%
% "Mandatory" (if missing will print the complete list of authors,
% including the \thanks symbols)
\authorref{Author1,~F., Author2,~S., and Author3,~T.}
%
% (Optional)
% \journalyear{2017}
% \journalvolume{V}
% \journalissue{N}
% \journalpages{xx--xx}
% \doi{xx.xxxx/xxxx.xx}

% Remaining Pages (Optional)
%
\authorshort{Author1 et al.} %or, e.g., \authorshort{Author1 et al}
\titleshort{Data-Driven Analysis of Form and Harmonic Structure in AI-Generated Popular Music}

%%%%%%%%%%%%%%%%%%%%%%%%%%%%%%%%%%%%%%%%%%%%%%%%%%%%%%%%%%%%%%%%%%%%%%%%%%%%%%%%
% Document Content
%%%%%%%%%%%%%%%%%%%%%%%%%%%%%%%%%%%%%%%%%%%%%%%%%%%%%%%%%%%%%%%%%%%%%%%%%%%%%%%%

\begin{document}

%%%%%%%%%%%%%%%%%%%%%%%%%%%%%%%%%%%%%%%%%%%%%%%%%%%%%%%%%%%%%%%%%%%%%%%%%%%%%%%%
% Abstract
%%%%%%%%%%%%%%%%%%%%%%%%%%%%%%%%%%%%%%%%%%%%%%%%%%%%%%%%%%%%%%%%%%%%%%%%%%%%%%%%

\twocolumn[{%
%
\maketitleblock
%

\begin{abstract}
We present a large-scale comparison of the formal and harmonic structures of music generated by users of two AI music generation platforms, Suno and Udio, against human-composed popular music. We address three questions: 1) How does AI-generated music's harmony relate to human popular music? 2) What structural similarities or differences exist? 3) What considerations arise when applying music information retrieval (MIR) tools to AI-generated music?
We analyze $20,002$ Suno tracks, $20,002$ Udio tracks, and $20,010$ human-composed recordings. Harmonic analysis uses a state-of-the-art conformer-based chord estimation model; structural analysis uses a joint beat-tracking and semantic section labeling model (All-In-One). We model chord progressions as trigrams and tetragrams, and musical form as n-grams of section labels, to quantify harmonic and formal diversity across the three collections.
Analysis of over $2.52$ million chord trigrams reveals that Suno concentrates its harmony on Western pop patterns such as \texttt{I-V-vi-IV}, while Udio exhibits more balanced chord distributions resembling human-composed music. Section-label analysis shows Suno tends toward simple verse-chorus repetitions, whereas Udio exhibits the widest structural variety. These findings suggest that Suno (versions 3.5 and below) converges on a narrow harmonic and formal space, whereas Udio exhibits tendencies closer to human-composed music. We release beat-level chord annotations, metadata, section boundaries, self-similarity matrices, audio links, and code under an MIT license.
\end{abstract}

%
\begin{keywords}
AI Music, Generative AI, Musical Form, Harmonic Structure, Automatic Chord Estimation, Exploratory Data Analysis
\end{keywords}
}]
\saythanks{}

%%%%%%%%%%%%%%%%%%%%%%%%%%%%%%%%%%%%%%%%%%%%%%%%%%%%%%%%%%%%%%%%%%%%%%%%%%%%%%%%
% Main Content Start
%%%%%%%%%%%%%%%%%%%%%%%%%%%%%%%%%%%%%%%%%%%%%%%%%%%%%%%%%%%%%%%%%%%%%%%%%%%%%%%%

\section{Introduction}
Commercial AI music generation platforms are becoming popular, and the resulting music is beginning to flood streaming services \citep{Sturm2024a}. Suno,\endnote{\url{https://suno.com}, last accessed Sep. 30 2025.}, launched online in late 2023, claimed in November 2024 to have a user base of 25 million people.\endnote{\url{http://bit.ly/4nnZA55}, last accessed Sep. 30 2025.} Udio,\endnote{\url{https://www.udio.com}, last accessed Sep. 30 2025.}
launched in the Spring of 2024, reported soon after that its users were generating ten songs a second.\endnote{\url{https://bit.ly/46QOGyy}, last accessed Sep. 30 2025.}
The summer of 2025 witnessed the appearance of an AI-generated music album by the fake ``indie rock band'' {\em The Velvet Sundown}, which has received millions of streams on Spotify.\endnote{\url{https://www.theguardian.com/technology/2025/jul/14/an-ai-generated-band-got-1m-plays-on-spotify-now-music-insiders-say-listeners-should-be-warned}, last accessed Sep. 30 2025.}
The French music streaming service Deezer announced in September 2025 that they are receiving $30,000$ ``fully AI-generated tracks daily''.\endnote{\url{https://www.musicbusinessworldwide.com/nearly-a-third-of-all-tracks-uploaded-to-deezer-are-now-fully-ai-generated-says-platform/}, last accessed Sep. 30 2025.}

This growing catalog of AI-generated music raises several interesting questions for musicology and presents challenges for music information retrieval (MIR). Given the rate at which this catalog is growing, how reliable and scalable are the computational methods of MIR for addressing these questions? How does the harmonic content of AI-generated music relate to the harmonic practices found in human-composed music? What structural similarities and differences are there between AI-generated and human-composed music?

In this article, we address these questions specifically in the context of AI-generated music, drawing on examples from Suno and Udio. We describe the creation of a large corpus of $60,014$ music recordings, balanced across three sources: $20,002$ songs from Suno users, $20,002$ songs from Udio users, and $20,010$ human-composed Western popular songs.

Recent advances in MIR have made corpus-scale harmonic and structural analysis of popular music increasingly viable. Even though the challenges of resolving chord and structural information across diverse musical types remain an open problem, we employ consonance-ACE \citep{poltronieri2025discord} for harmonic analysis and All-In-One \citep{taejun2023allinone} for structural analysis, two recent implementations that bring the analysis significantly closer to musicological utility. The results of our harmonic analyses reveal that music from Suno exhibits strong stylistic characteristics of Western pop chord progressions, whereas music from Udio exhibits more balanced chord-quality distributions, resembling those found in human-composed music. The results of our structural analyses show that music from Suno tends toward simple AB (verse-chorus) repetition, whereas music from Udio exhibits a wider range of structural variations.

In summary, our contributions are:
(i) a large comparative analysis of harmonic and formal properties in human-made and AI-generated music using a consistent pipeline,
(ii) publicly released annotations for $60,014$ tracks, including segmentations, beat-level chord estimates, style tags \citep{oramas2018multimodal}, and audio metadata, and (iii) open-sourced code to reproduce the experiments.\endnote{\url{https://anonymous.4open.science/r/Aimir_Structure_Analysis}}
% The remainder is organized as follows. The next section reviews related work. Section~\ref{sec:corpora} describes the corpora used in the study. Section~\ref{sec:method} motivates and describes our pipeline for harmonic and formal analysis. Section~\ref{sec:trigram_comparisons} presents our results. Finally, Sec.~\ref{sec:discussion} discusses implications of our work, its limitations, and future directions.

% -----------------------------------------------------------------
\section{Background}
\label{sec:background}

\rev{MIR provides computational methods for analyzing musical structure and harmony \citep{Muller2015a}. This section reviews the strands of that work the present study builds on: AI music generation and the evaluation of its output, the two methods we apply at corpus scale, music structure analysis and automatic chord estimation, and prior corpus-scale harmonic analyses of popular music.}

\subsection{AI Music Generation}
Recent open-source work on text-conditioned music and audio synthesis includes models such as MusicGen \citep{copet2023simple}, MusicLM \citep{agostinelli2023musiclm}, and Stable Audio Open \citep{evans2025stable}, representing the state of the art in publicly available generative audio research. However, the focus of this paper is not on the generative models themselves, but on the musical properties of their output, analyzed at the corpus scale in comparison with human-produced music.

\citet{Sturm2016c} evaluates \textit{folk-rnn}, a generative system trained on a large corpus of folk tunes in ABC notation, examining outputs in relation to musical practice rather than only automatic metrics, with a concrete evaluation pipeline, corpus statistics analysis, and focused musical inspection to locate where model outputs resemble and where they diverge from human repertoire. \citet{Ariza2009a} shows that interrogator-style authorship discrimination tests only establish that a digital music artifact sounds plausible to listeners, which is insufficient to characterize the musical quality produced by an AI system.

\citet{dervakos2021heuristics} propose a framework for evaluating symbolic AI-generated music using tools from music theory, defining heuristic metrics grounded in tonic coordinate systems and tone networks, such as the circle of fifths, to quantify musicality without relying on user surveys or real data.

\citet{wu2020jazz} analyze the outputs of the Jazz Transformer using objective metrics, including pitch-class histogram entropy, chord-progression irregularity, groove-pattern similarity, and structureness indicators, demonstrating that systematic shortcomings in AI-composed music are not fully captured by listening tests alone. Notably, their definition of chord progression irregularity as the proportion of unique chord trigrams in a piece establishes a precedent for trigram-based harmonic analysis of AI-generated music.

Some studies find biased perception of music when works are labeled as AI \citep{Ragot2020t,Shank2022q,Deguernel2023a}, others report null effects \citep{Pasquier2016a,Zlatkov2023a}. Together, these motivate a careful, data-driven comparison between music made with current AI platforms and music created by humans.

\subsection{Music Structure Analysis}
\label{sec:msa}
Music Structure Analysis (MSA) is the task of segmenting an audio recording into contiguous, non-overlapping sections and assigning labels that reflect musical function or similarity \citep{nieto2020audio}. Computational approaches exploit four main principles: homogeneity (sections internally consistent with respect to some musical attribute), repetition (detecting recurrent sequences), novelty (identifying boundary points where musical content changes abruptly), and regularity (statistical properties of segment durations). Classical methods relied on self-similarity matrices and novelty curves derived from hand-crafted features such as chromagrams and MFCCs \citep{foote2000automatic}, repetition and recurrence-based representations \citep{levy2008structural, serra2012unsupervised}, Laplacian spectral decomposition \citep{mcfee2014learning}, and barwise self-similarity with dynamic programming \citep{Marmoret_2023}. Despite these advances, MSA remains an open challenge due to the subjectivity of structural annotations: the best reported structure labeling scores reach approximately $74\%$ on the Harmonix Set \citep{taejun2023allinone}, well below the estimated human inter-rater agreement of approximately $90\%$ \citep{nieto2020audio}.

More recent approaches have improved structural analysis through deep learning, including few-shot embeddings for Laplacian decomposition \citep{salamon2021deep}, unsupervised contrastive multi-level representations \citep{buisson2022learning}, self-supervised novelty detection \citep{peeters2023self}, hierarchical spectral-temporal Transformers for direct semantic labeling \citep{wang2022catch}, and graph attention networks for pairwise structure prediction \citep{buisson2024using}.

For functional structure extraction, we employ the All-In-One model \citep{taejun2023allinone}, which jointly performs beat tracking, downbeat detection, segment boundary detection, and semantic section labeling (verse, chorus, bridge, etc.) in a single architecture. \rev{We adopt it for two reasons: it directly outputs the section labels our analysis requires, and it reaches state-of-the-art performance on the Harmonix Set across all four tasks.}

\subsection{Automatic Chord Estimation}
\label{sec:ace}
Automatic Chord Estimation (ACE) is a central task in MIR with applications in music performance, transcription, and formal analysis research. ACE has evolved from early template and chroma-based systems with probabilistic sequence models, such as Hidden Markov Models, Conditional Random Fields, or Deep Belief Networks \citep{mauch2009simultaneous, ni2012end, humphrey2015four}, to more recent deep learning approaches that learn end-to-end representations from spectral features, increasingly targeting large chord vocabularies that include sevenths, extensions, and inversions.

Large-vocabulary ACE has advanced through several complementary strategies. \citet{mcfee2017structured} decouple pitch-class detection from label assignment in a 170-class framework; \citet{jiang2019large} decompose chord labels into independent musical components (root, quality, bass, extensions), generalizing to 301+ classes; and \citet{park2019bi} propose BTC, a bi-directional Transformer achieving competitive large-vocabulary performance with publicly available code. Most recently, \citet{poltronieri2025discord} propose consonance-ACE, a Conformer-based model whose decomposed decoder reconstructs chord labels without a fixed vocabulary, paired with consonance-based label smoothing, outperforming BTC on both standard and non-binary evaluation metrics on RWC Pop and USPop. \citet{waseem2025chordformer} further extends this line with ChordFormer for large-vocabulary recognition.

Despite these advances, implementations remain limited in coverage. Open-source libraries such as Madmom \citep{bock2016madmom} and librosa \citep{mcfee2025librosa} \rev{implement} baseline chord estimation algorithms, most of which remain fundamentally limited to detecting major and minor triads, struggling with tetrads, extensions, suspensions, and inversions. Earlier efforts to expand chord vocabularies, such as \textit{Chordino} \citep{mauch2010approximate}, \rev{are no longer maintained}. Furthermore, analyses relying on chord-class aggregation can misrepresent stylistic progressions: \citet{collins2025recording} demonstrates that \textit{Chordino}-based histograms failed to capture modal-interchange progressions while detecting only broad palette differences. Multi-source approaches such as DECIBEL \citep{odekerken2020decibel}, which fuse audio features with crowd-sourced MIDI annotations and guitar tablature, partially address vocabulary restrictions, yielding improvements of 0.5--13.6 percentage points over audio-only baselines on $180$ Beatles songs and $20$ Queen songs from the Isophonics chord-annotation corpus\endnote{\url{http://isophonics.net/content/reference-annotations}, last accessed Sep. 30 2025.}, though major and minor accuracy typically remains between $75\%$ and $87\%$, a ceiling attributed to annotation subjectivity.

Roman Numeral Analysis (RNA) encodes chords relative to a tonal center. Benchmarks on the PARC dataset \citep{poppe2025polyphonic} confirm that \rev{models predicting Roman numerals directly from audio} achieve modest accuracy, particularly for complex progressions. \rev{We therefore derive Roman numerals in two stages rather than directly from audio: ACE recognizes the chords, and a symbolic step assigns each its key-relative function through global key estimation (Section~\ref{sec:ace_pipeline}).}

Given these considerations, this study adopts consonance-ACE \citep{poltronieri2025discord} for three reasons: (1) its decomposed decoder is not constrained to a fixed vocabulary, accommodating harmonic diversity beyond major/minor triads; (2) consonance-based label smoothing produces errors that fall between harmonically related chords rather than random ones; and (3) its demonstrated superiority over BTC on standard benchmarks, combined with publicly available code, provides a reproducible foundation for large-scale analysis. The model is described in detail in Section~\ref{sec:ace_pipeline}.

\subsection{Harmonic Analysis of Corpora}
Several annotated corpora and computational analyses have contributed to the study of harmony and musical structure in popular music. \rev{The hand-annotated rock corpus of \citet{temperley2013statistical} supports statistical investigation of chords, scale degrees, and key-finding models.} Expert ground-truth chord annotations for charting songs support research in automatic chord recognition and music analysis~\citep{burgoyne2011expert}. Computational analysis of Billboard chart data has been used to study the evolution of harmony, timbre, and style in popular music from 1960 to 2010~\citep{mauch2015evolution}. Long-term trends in harmonic and timbral features of contemporary Western popular music have also been examined to document changes in musical diversity~\citep{serra2012measuring}. Automated approaches to chord and key estimation provide frameworks for the simultaneous identification of chords and musical context from audio data~\citep{mauch2009simultaneous}. However, there is no comprehensive corpus-level comparison of harmonic and structural properties between current AI platforms and human music.

% ------------------------------------------------------------------
\section{Corpora}
\label{sec:corpora}

The corpora are formed as a collection of three audio datasets (summarized in Table \ref{tab:datasets}) containing audio files from \textit{Suno} and \textit{Udio} (AI-generated) and human-generated music from the Lastfm catalog. The audio tracks originate from \citep{cros2025ai}, where web scraping was performed for the \textit{Suno} and \textit{Udio} datasets between May and October 2024. \textit{Suno} data was collected from its frequently updated ``New Songs" playlist \endnote{\url{https://suno.com/playlist/cc14084a-2622-4c4b-8258-1f6b4b4f54b3}, last accessed Sep. 30, 2025.}, queried every two hours to capture its daily influx of $50$ new tracks. \textit{Udio} data was collected using the platform's search functionality, retrieving songs based on their popularity by tracking those with the highest number of plays in the 24 hours before each query \citep{cros2025ai}. Both datasets include metadata and audio files. The audio file folder contained more than $200,000$ files, from which we randomly selected approximately $20,000$ audio files from each dataset for analysis. The corpora uses the Million Song Dataset (MSD) \citep{Bertin-Mahieux2011}\endnote{\url{https://github.com/renesemela/lastfm-dataset-2020/raw/master/datasets/lastfm_dataset_2020/lastfm_dataset_2020.db}, last accessed Sep. 30 2025.} to confirm the audio files are human-created music, leveraging Lastfm's associated YouTube URLs to download around $200,000$ unique tracks, where we randomly chose $20,010$.\endnote{We used \texttt{yt-dlp} \url{https://github.com/yt-dlp/yt-dlp}, last accessed Sep. 30 2025.} We thus refer to this collection as ``\textit{Lastfm}''.


\begin{table}[htpb]
\centering
\footnotesize
\begin{tabular}{llll}
\toprule
\textbf{Dataset} & \textbf{\# JSON} & \textbf{Duration} & \textbf{Total} \\
& \textbf{\# Files} & \textbf{Mean $\pm$ S.D.} & \textbf{Duration} \\
\midrule
    Lastfm & 20,010 & 261.3 $\pm$ 182.1s & 1452h 6m 56s \\
    Suno & 20,002 & 175.2 $\pm$ 54.6s & 1084h 47m 9s \\
    Udio & 20,002 & 143.2 $\pm$ 98.6s & 801h 47m 36s \\
\bottomrule
\end{tabular}
\caption{The number of files is formed by a parallelized process with workers extracting form and chords data from a randomized folder of more than $100,000$ audio files, with a limit of $20,000$ JSON files}
\label{tab:datasets}
\end{table}

We report user engagement metrics (likes and plays) across the Suno and Udio collections. Table~\ref{tab:engagement} presents summary statistics, and Figure~\ref{fig:engagement} displays their distributions on a logarithmic scale. Both platforms share identical median values (1 like and 4 plays), with means driven upward by a small number of high-engagement outliers. The 90th percentile reaches only $2$ likes and $26$ plays for Suno, and $3$ likes and $32$ plays for Udio, confirming that the vast majority of tracks in both collections received minimal engagement.

\begin{table}[ht]
\centering
\footnotesize
\begin{tabular}{llrrrrr}
\toprule
\textbf{Metric} & \textbf{Collection} & \textbf{Mean} & \textbf{Med.} & \textbf{P90} & \textbf{P95} & \textbf{Max} \\
\midrule
Likes & Suno  &   5.25 & 1 &  2 &   4 &  34{,}921 \\
      & Udio  &   1.83 & 1 &  3 &   8 &     898   \\
\midrule
Plays & Suno  & 204.88 & 4 & 26 &  50 & 445{,}490 \\
      & Udio  &  44.10 & 4 & 32 &  58 &  48{,}940 \\
\bottomrule
\end{tabular}
\caption{Summary statistics of likes and plays for the Suno ($n=19{,}972$) and Udio ($n=19{,}992$) collections. P90 and P95 denote the 90th and 95th percentiles.}
\label{tab:engagement}
\end{table}

\begin{figure}[ht]
    \centering
    \includegraphics[width=\linewidth]{figures/engagement_boxplot.png}
    \caption{Box plots of likes and plays for the Suno and Udio collections on a logarithmic scale. Both platforms share a median of 1 like and 4 plays.}
    \label{fig:engagement}
\end{figure}

% ------------------------------------------------------------------
\section{Method}
\label{sec:method}
This section describes the pipeline used to extract harmonic and structural information from the $60,014$ audio recordings in our corpus, and the derived representations for comparative analysis.

\subsection{Automatic Chord Estimation}
\label{sec:ace_pipeline}
Chord estimation is performed using consonance-ACE \citep{poltronieri2025discord}, which independently predicts root, bass, and pitch-class activations and reconstructs chord labels symbolically through a decomposed decoder. The model operates on CQT spectrograms computed over 20-second chunks and outputs frame-level chord predictions with no constraint to a fixed vocabulary. We use the publicly available pretrained model without modification. Chord segments shorter than $0.5$ seconds are merged into adjacent segments, and consecutive frames with identical chord labels are collapsed into continuous intervals.

The output of the chord estimator consists of absolute chord labels (e.g., \texttt{D:min7}, \texttt{G:7}, \texttt{C:maj7}) aligned with a time onset-offset columns as shown in Table \ref{tab:ace_output}. To enable key-relative analysis, we convert these to functional harmony representations (e.g., \texttt{ii}, \texttt{V}, \texttt{I}). \rev{For each track, we first estimate a single global key from the chord sequence by duration-weighted diatonic template matching, scoring each of the 24 major and minor keys by the time spent on chords whose root and quality fit its scale degrees.} Each absolute chord label is then mapped to its corresponding Roman numeral degree relative to the estimated key using the \texttt{music21} library \citep{cuthbert2010music21}. Consecutive identical Roman numerals are de-duplicated before downstream n-gram analysis. The resulting chord intervals are subsequently aligned to beat positions estimated by the All-In-One model (Section~\ref{sec:form_pipeline}), yielding the beat-level chord annotations included in the released dataset.

\begin{table}[ht]
\centering
\footnotesize
\begin{tabular}{rrl}
\toprule
\textbf{Start (s)} & \textbf{End (s)} & \textbf{Chord} \\
\midrule
0.000 & 0.951 & N \\
0.951 & 2.900 & G:min \\
2.900 & 4.501 & F:maj \\
4.501 & 7.471 & D\#:maj \\
7.471 & 9.118 & G:min \\
9.118 & 11.346 & F:maj \\
11.346 & 12.691 & C:min \\
12.691 & 13.805 & D\#:maj \\
13.805 & 15.360 & G:min \\
15.360 & 16.914 & F:maj \\
16.914 & 18.840 & D\#:maj6 \\
18.840 & 21.787 & G:min \\
21.787 & 23.179 & F:maj \\
23.179 & 24.617 & D\#:maj \\
24.617 & 25.499 & G:min \\
25.499 & 26.241 & D\#:maj \\
\bottomrule
\end{tabular}
\caption{Example output of consonance-ACE. Each row contains the onset time, offset time, and estimated chord label. N denotes no chord detected. \rev{The model predicts a root pitch class and quality, labeling roots with sharps by default and assuming enharmonic equivalence.}}
\label{tab:ace_output}
\end{table}

\subsection{Musical Form Estimation}
\label{sec:form_pipeline}
Musical structure is extracted using the All-In-One model \citep{taejun2023allinone}, which jointly performs four tasks in a single architecture: beat tracking (locating regular pulse positions), downbeat detection (identifying the first beat of each bar), segment boundary detection (finding time points where sections change), and functional structure labeling (assigning semantic labels such as verse, chorus, bridge, intro, outro, or solo to each segment).

The model first separates the input audio into four stems (drums, bass, vocals, other) using Hybrid Transformer Demucs \citep{rouard2023hybrid} and processes the resulting spectrograms through a stack of transformer modules combining dilated and neighborhood attention. Four output heads produce frame-level predictions for each task. The model achieves state-of-the-art performance on the Harmonix Set ($F1 = 0.958$ for beats, $F1 = 0.915$ for downbeats, $HR.5F = 0.660$ for segmentation, $PWF = 0.738$ for structure labeling).

We apply the publicly available pretrained model without modification. For each track, the model outputs a sequence of labeled sections with their onset and offset times, stored as \texttt{\_musical\_form.lab} files. Consecutive identical section labels are subsequently merged before downstream analysis.

\subsection{Trigram Dataset Generation}
\label{sec:trigram_generation}
To systematically analyze harmonic tendencies across collections, we constructed a \textit{trigram dataset}, a corpus of all consecutive three-chord sequences extracted from each collection's data.
For example, the progression $Dm7$, $Ab7$, $Gm7$, $C\#ø7$, \ldots\ yields $[Dm7, Ab7, Gm7]$, $[Ab7, Gm7, C\#ø7]$, and so on.
The process is as follows:

\begin{enumerate}
    \item \textbf{Normalization:} Absolute chord labels are first converted to Roman numerals, and consecutive identical Roman numerals are de-duplicated, as described in Section~\ref{sec:ace_pipeline}. The example above uses absolute chord labels for illustrative clarity. Sequences containing ``no-chord'' labels were not considered for trigrams.

    \item \textbf{Trigram Generation:}
    For each song, we extracted overlapping chord trigrams from the Roman numeral sequence using a sliding window approach:
    \[
        (C_{i},\ C_{i+1},\ C_{i+2}), \quad i = 1, \ldots, N-2
    \]
    where $N$ is the number of chord events in the song after de-duplication, and $C_i$ denotes the $i^{th}$ Roman numeral. All trigrams were retained, including repetitions and sequences spanning sections or boundaries, unless a song or section was shorter than three chords.
    
\end{enumerate}

\subsection{Tetragram Dataset Generation}
\label{sec:tetragram_generation}

In parallel with the trigram dataset (Section~\ref{sec:trigram_generation}), we construct a \textit{tetragram dataset}, a corpus of all consecutive four-chord sequences from each collection, applying the same sliding-window principle to the Roman numeral sequences produced by the ACE pipeline (Section~\ref{sec:ace_pipeline}).

For each song, consecutive identical Roman numerals are first de-duplicated, as in Section~\ref{sec:ace_pipeline}, and ``no-chord'' labels are excluded. A sliding window of four chords $(C_{i},\,C_{i+1},\,C_{i+2},\,C_{i+3})$, $i = 1, \ldots, N-3$, is then extracted, where $N$ is the number of chord events after de-duplication and $C_i$ denotes the $i^{\text{th}}$ Roman numeral.

A window is retained only if at least three of its four chords are distinct, compared case-insensitively so that a degree and its mode (e.g.\ \texttt{IV} and \texttt{iv}) count as one; songs shorter than four chords yield no tetragrams. We refer to these four-chord windows as \textbf{tetragrams}. They are counted per collection, revealing the most common four-chord progressions and their cross-collection frequency distributions.

\subsection{Axis of Awesome Chord Progression Analysis}
\label{sec:axis}

We measure how often songs in each collection contain chord progressions matching the Axis of Awesome \citep{richards2017tonal}, using the Roman numeral sequences produced by the ACE pipeline. %(Section~\ref{sec:ace_pipeline}). 
\rev{Four-chord sequences, which we refer to as \textbf{tetragrams}, are sliding four-chord windows over the deduplicated Roman numeral sequence, as in Section~\ref{sec:tetragram_generation}; here every window enters the analysis, without the retention rule used to build the tetragram dataset.}

The Axis of Awesome is associated with three four-chord sets: the major form \{I, V, vi, IV\}, the minor form \{i, $\flat$VI, III, $\flat$VII\}, and a minor variant \{i, $\flat$VI, $\flat$III, $\flat$VII\}. \rev{We report two measures of how a song uses these sets. The first, axis-family membership, counts a four-chord window when its chords form one of the sets in any of the $24$ orderings ($72$ patterns in total), and measures whether a track draws on axis-family pitch material. The second, the literal Axis loop, counts a window only when it is one of the $12$ rotations that keep the chords in their canonical cyclic order, and measures the recurring progression itself. This follows \citet{richards2017tonal}: the cycle strongly tends to appear in its given order. Rotation equivalence also absorbs relative-major/minor ambiguity in the global key estimate from which the Roman numerals are derived (Section~\ref{sec:ace_pipeline}): the four chords C, G, Am, F are labeled \texttt{I-V-vi-IV} when the key is read as C~major but \texttt{III-bVII-i-bVI}, the minor-form set begun on a different chord, when it is read as the relative A~minor. Treating all three sets and their rotations as one loop recovers the progression in both readings, so a track is not counted as non-axis merely because its key estimate landed on the relative key.} From here, we will refer to this pattern of chord progressions as Axis.

All Roman numeral labels are first normalized enharmonically via a 12-tone semitone mapping (e.g., \texttt{\#V}~$\rightarrow$~$\flat$\texttt{VI}), and consecutive identical labels are removed. A sliding window of four chords is then passed over each collapsed sequence, and every four-chord window $(a, b, c, d)$ is \rev{tested for membership in both pattern sets: the $72$ axis-family patterns and the $12$ rotation-only patterns}.

\rev{The \textbf{tetragram axis rate} for a song is the proportion of its four-chord windows that match, expressed as a percentage; we compute it separately under the axis-family and the rotation-only measures.}

\begin{figure*}[!t]
    \centering
    \includegraphics[width=\linewidth]{figures/genres.png}
    \caption{Top 30 most common trigrams found in Lastfm, Suno, and Udio datasets. The y-axis shows the relative frequency of each trigram in the collection. The distribution of styles differs across collections.}
    \label{fig:top_30_trigrams}
\end{figure*}

% ---------------------------------------------------------------
\section{Comparing Patterns in Form and Harmony}
\label{sec:trigram_comparisons}
The analyses proceed from the broadest view of the harmonic palette to increasingly specific patterns of chord succession and formal organization.

% Trigram tendencies
% Overal group of trigrams per collection
% SSM with BPE
% Chord detection accuracy comparison

\subsection{Trigram Tendencies Across Collections}
\label{sec:trigram_tendencies}

Figure~\ref{fig:top_30_trigrams} shows the thirty most prevalent trigram patterns in the three members of the collection (Lastfm, Suno, and Udio), distributed over the ten most common styles found by extracting the metadata descriptors of the songs. The label ``Others" collects several styles in one group. Both Udio and Suno show a strong tendency toward Pop, Rock, Metal, and Electronic.


% Note: Describe better the style detection in the three collections. 
While we acknowledge that the style tags associated with tracks in our dataset are inherently flawed and subject to user bias (users may arbitrarily assign genres, and the platforms' own tagging systems lack rigorous musicological grounding), we use these classifications as they reflect the styles users most frequently specify when generating music. As demonstrated by \rev{\citet{casini2025data}}, these same style descriptors (Pop, Rock, Electronic, Hip-hop, etc.) are among the most prevalent tags in the dataset.

The collection-level differences visible in Figure~\ref{fig:top_30_trigrams} are quantified in detail in Sections~\ref{sec:suno_tendencies}--\ref{sec:axis_tendencies}.

In Figure~\ref{fig:vocabulary_distribution}, an equivalent distribution of chord qualities is seen in all three collections. The major chord is the most common, followed by minor, power chords, suspended fourths, and dominant sevenths. As an observation, dominant chords with a sharp ninth or flat thirteenth, also altered dominant, are commonly used in minor and major modes, as a modal interchange in major tonalities. The dominant sharp eleventh is not common; the once-ubiquitous dominant substitution used in jazz, bossa-nova, or soul has become an endangered harmonic particularity in all three collections. Power chord, as a neutral quality, is used with mixed frequency across major and minor tonalities.

\begin{figure}[ht]
    \centering
    \includegraphics[width=\linewidth]{figures/vocabulary_ace.png}
    \caption{Vocabulary distribution of chord qualities over the collections. The y-axis indicates the total number of occurrences for each chord quality in each collection, scaled by $\times 10^{5}$.}
    \label{fig:vocabulary_distribution}
\end{figure}

Figure~\ref{fig:tonalities} shows a distribution of the tonalities found in all collections. The quantification is sorted to Lastfm occurrence. The proportions show similarities on AI-generated music platforms; nevertheless, Suno shows stronger tendencies to produce C major, D major, G major, and B major, compared to Lastfm and Udio. Udio has a prominent presence of A minor, G minor, and Eb minor. The distribution of tonalities is consistent with the cycle of the fifth.

\begin{figure}[ht]
    \centering
    \includegraphics[width=\linewidth]{figures/tonalities_per_collection_ace.png}
    \caption{Distribution of most common tonalities in all three collections.}
    \label{fig:tonalities}
\end{figure}

\rev{Figure~\ref{fig:trigrams_per_style} shows the harmonic palette of each collection across the eight most common styles (Rock, Electronic, Pop, Jazz, Country, Hip-Hop, Soul, and Blues): for every style it reports the most frequent chord trigrams and how each collection distributes its usage across them. To compare the palettes by shape, each collection is normalized to its own maximum within a style, so the figure describes which progressions a collection relies on, not how often it uses them. Read this way, the three collections share a common canonical core but differ in emphasis. How frequently these progressions occur, and how far Suno over-represents them, is quantified in Section~\ref{sec:suno_tendencies} (Figure~\ref{fig:pop_prevalence}, Table~\ref{tab:suno_enrichment}).}

\begin{figure*}[!htbp]
    \centering
    \includegraphics[width=\linewidth]{figures/radial_chart_trigrams_ace.png}
    \caption{\rev{Per-style harmonic palette. For each style, the axes are the most frequent chord trigrams (ordered by Lastfm), and each collection is scaled to its own maximum within that style; the figure therefore shows which progressions each collection draws on, not how often. Polygon size carries no magnitude. Absolute frequencies appear in Figure~\ref{fig:pop_prevalence} and Table~\ref{tab:suno_enrichment}.}}
    \label{fig:trigrams_per_style}
\end{figure*}

\subsection{Harmonic Tendencies}
\label{sec:suno_tendencies}
\rev{Figure~\ref{fig:pop_prevalence} shows the relative frequency of canonical pop/rock chord trigrams, defined as permutations of \{I/i, IV/iv, V/v, vi/VI/$\flat$VI\}, in each collection.} The figure demonstrates that Suno tends to favor chord progressions strongly associated with the Western pop idiom, particularly those prevalent since the early 2000s. For every progression tested, Suno's relative frequency is significantly greater than that observed in the other collections ($p < 0.001$ for all comparisons). As shown in ~\citep{temperley2013statistical}, the centrality of progressions such as \texttt{I-V-vi-IV}, \texttt{IV-I-V}, and related idioms is a defining feature of Western pop/rock music; they report that the most frequent non-tonic chord roots are \texttt{IV}, \texttt{V}, \texttt{bVII}, and \texttt{vi}, with progressions among these chords forming the core stylistic harmonic vocabulary. Our results both replicate and extend these findings, revealing the extent to which such progressions are even more prevalent in AI-generated music from Suno.

To quantify the tendencies shown in Figure~\ref{fig:trigrams_per_style} for each canonical progression in Suno relative to the reference collections (Lastfm and Udio), we calculated a \emph{risk ratio (RR)} for every trigram.  Raw counts are misleading because Suno contains far more trigrams than the reference corpus, so the RR normalizes by each collection's total:

\[
\operatorname{RR}(t) \;=\;
\frac{N_{\mathrm{Suno},t}/N_{\mathrm{Suno}}}
     {N_{\mathrm{Ref},t}/N_{\mathrm{Ref}}}
\]

\noindent
where $N_{\mathrm{Suno},t}$ and $N_{\mathrm{Ref},t}$ are the counts of trigram~$t$ in \textit{Suno} and in the reference corpus, respectively, while $N_{\mathrm{Suno}}$ and $N_{\mathrm{Ref}}$ denote the total numbers of trigrams in those two corpora.

\smallskip
To test whether the difference is greater than expected by chance, we build a $2\times2$ contingency table (\emph{Suno vs.\ Reference}\,$\times$\,$t$ vs.\ ``all other trigrams'') and apply Fisher's exact test. All raw $p$-values are corrected for multiple comparisons with the Benjamini--Hochberg false-discovery-rate procedure \citep{benjamini1995controlling}; the resulting $p$-value is used to filter and choose the trigrams that are shown in Table \ref{tab:suno_enrichment}:

\begin{itemize}
  \item at least $30$ occurrences in Suno and $10$ in the reference,
  \item risk ratio $RR\ge2$,
  \item adjusted $p-$value $ < 0.01$.
\end{itemize}

\begin{center}
\begin{tabular}{lcc}
                 & Target trigram & Other trigrams \\\hline
\textbf{Suno}    & $a$            & $b$ \\
\textbf{Reference} & $c$            & $d$
\end{tabular}
\end{center}

\noindent where \(a=N_{\text{Suno},t}\), \(b=N_{\text{Suno}}-a\), \(c=N_{\text{Ref},t}\), and \(d=N_{\text{Ref}}-c\).
A risk ratio $RR_t > 1$ together with $p< 0.01$ indicates that trigram $t$ is significantly more common in Suno than in the reference corpus, as is the case for all the trigrams shown.

\begin{figure}[!ht]
    \centering
    \includegraphics[width=\linewidth]{figures/pop_prevalence_suno.png}
    \caption{\rev{Relative frequency of canonical pop/rock trigrams in \textit{Suno} compared with \textit{Lastfm} and \textit{Udio}. Bars show the relative frequency of each progression within a collection, and all differences between Suno and the reference collections are significant at \(p<0.01\).}}
    \label{fig:pop_prevalence}
\end{figure}

Table~\ref{tab:suno_enrichment} shows a representative selection of the $118$ Suno-enriched trigrams passing all filters: Panel~A lists the top~15 by Suno count of the 53~canonical progressions; % (Section~\ref{sec:suno_tendencies}); 
Panel~B lists the top~10 non-canonical trigrams by RR from the remaining 65. The columns provide the raw counts in Suno ($C_\text{Suno}$), Udio ($C_\text{Udio}$), and Lastfm ($C_\text{Last}$), their percentage of the total vocabulary, and the RR relative to Lastfm.

Across all canonical progressions in Panel~A, RRs range from $2.1\times$ to $3.3\times$, with the major-mode axis chains (e.g., \texttt{I-V-vi} at $3.2\times$, \texttt{vi-IV-I} at $2.7\times$) among the most enriched. Every canonical pattern is significantly enriched ($\texttt{p-adj} < 10^{-16}$).

\begin{table}[!ht]
\centering
\tiny
\setlength{\tabcolsep}{4pt}
\begin{tabular}{@{}lrrrrrrc@{}}
\toprule
Trigram & $C_{\text{Suno}}$ & $C_{\text{Udio}}$ & $C_{\text{Last}}$ & \% Suno & \% Udio & \% Last & RR \\
\midrule
\multicolumn{8}{@{}l}{\textit{Panel A: Top 15 canonical enriched trigrams}} \\
\midrule
vi, IV, I    & 34095 & 6442 & 8659 & 4.326 & 1.770 & 1.596 & 2.71 \\
V, vi, IV    & 29383 & 6311 & 7311 & 3.728 & 1.734 & 1.347 & 2.77 \\
I, V, vi     & 22803 & 3967 & 4948 & 2.893 & 1.090 & 0.912 & 3.17 \\
V, i, bVI    & 4823  & 1221 & 1077 & 0.612 & 0.336 & 0.198 & 3.08 \\
v, i, bVI    & 4610  &   863  & 1022 & 0.585 & 0.237 & 0.188 & 3.11 \\
iv, V, i     & 4499  & 1373 & 1029 & 0.571 & 0.377 & 0.190 & 3.01 \\
bVI, V, i    & 3508  & 1092 & 1074 & 0.445 & 0.300 & 0.198 & 2.25 \\
iv, v, i     & 3287  &   653  &   739  & 0.417 & 0.179 & 0.136 & 3.06 \\
V, i, iv     & 3273  & 1280 &   950  & 0.415 & 0.352 & 0.175 & 2.37 \\
bVI, v, i    & 3117  &   618  &   658  & 0.396 & 0.170 & 0.121 & 3.26 \\
i, v, bVI    & 2563  &   594  &   650  & 0.325 & 0.163 & 0.120 & 2.72 \\
v, i, iv     & 2553  &   650  &   762  & 0.324 & 0.179 & 0.140 & 2.31 \\
i, iv, V     & 2128  &   803  &   576  & 0.270 & 0.221 & 0.106 & 2.54 \\
i, bVI, v    & 2081  &   554  &   555  & 0.264 & 0.152 & 0.102 & 2.58 \\
i, iv, v     & 1929  &   581  &   622  & 0.245 & 0.160 & 0.115 & 2.14 \\
\midrule
\multicolumn{8}{@{}l}{\textit{Panel B: Top 10 non-canonical enriched trigrams (by RR)}} \\
\midrule
iv, ii, v    &    93 &  17 &  10 & 0.012 & 0.005 & 0.002 & 6.40 \\
iii, vi, ii  & 1485 & 236 & 181 & 0.188 & 0.065 & 0.033 & 5.65 \\
bVI, iv, ii  &    74 &  25 &  10 & 0.009 & 0.007 & 0.002 & 5.10 \\
II, v, III   &   245 &  30 &  35 & 0.031 & 0.008 & 0.006 & 4.82 \\
IV, iii, vi  & 1390 & 276 & 199 & 0.176 & 0.076 & 0.037 & 4.81 \\
V, biv, vi   &   198 &  33 &  30 & 0.025 & 0.009 & 0.006 & 4.54 \\
IV, biv, vi  &   131 &  29 &  21 & 0.017 & 0.008 & 0.004 & 4.30 \\
biv, vi, IV  &   195 &  39 &  32 & 0.025 & 0.011 & 0.006 & 4.20 \\
iii, IV, v   &   143 &  26 &  24 & 0.018 & 0.007 & 0.004 & 4.10 \\
iii, ii, v   &    67 &  19 &  12 & 0.009 & 0.005 & 0.002 & 3.84 \\
\bottomrule
\end{tabular}
\caption{Suno-enriched trigrams. Panel~A lists the most prevalent canonical progressions, defined as any permutation of chords drawn exclusively from \{I/i, IV/iv, V/v, vi/VI/bVI\}. Panel~B lists the highest RR non-canonical trigrams. $C$ = raw occurrence count; RR = risk ratio (Suno\,/\,Lastfm). All rows are statistically significant (FDR-adjusted $p < 0.001$).}
\label{tab:suno_enrichment}
\end{table}

\subsection{N-gram Analysis of Section-Label Sequences}
\label{sec:form_ngram}

% To explore common patterns in how songs are structured, we used Byte Pair Encoding (BPE) on sequences of labeled segments. 
We analyze sequences of section labels to explore patterns in song structure, using the same n-gram construction applied to chords (Sections~\ref{sec:trigram_generation} and~\ref{sec:tetragram_generation}).
For each track, the section sequence is read from the All-In-One output, consecutive identical sections are collapsed, and trigrams and tetragrams of consecutive section labels are extracted with a sliding window (e.g.\ \texttt{verse·chorus·verse}).

This yields $104{,}865$ section-label trigrams in Lastfm, $71{,}870$ in Suno, and $56{,}232$ in Udio ($232{,}967$ in total), together with the corresponding tetragrams ($86{,}086$, $53{,}704$, and $42{,}751$, respectively). Their cross-collection frequencies reveal the structural motifs that reflect how songs are typically organized.
Figure~\ref{fig:musical_form_heatmap} shows the most frequent section-label tetragrams; the results are discussed in Section~\ref{sec:disc_structure}.
\begin{figure}[!t]
    \centering
    \includegraphics[width=\linewidth]{figures/tetragram_musical_form.png}
    \caption{Section-label tetragram heatmap across collections. Each row is a tetragram of four consecutive section labels, extracted after collapsing repeated adjacent sections. Cell shading encodes relative frequency within each collection; integer annotations report absolute counts. Tetragrams are sorted by descending Lastfm frequency.}
    \label{fig:musical_form_heatmap}
\end{figure}



% Add SSM visualization of lastfm, suno and udio

\begin{figure*}[!ht]
    \centering
    \includegraphics[width=\linewidth]{figures/radial_musical_form.png}
    \caption{Normalized musical form trigram profiles across eight human-annotated genre tags per collection. Each axis represents a trigram of three consecutive section labels (e.g., chorus-verse-chorus), extracted after collapsing repeated adjacent sections. For each genre, the top trigrams are selected by Lastfm frequency, and values are normalized to the maximum count per collection within each genre.}
    \label{fig:musical_form}
\end{figure*}

\begin{figure*}[!ht]
    \centering
    \includegraphics[width=\linewidth]{figures/mergedPlot.pdf}
    \caption{Most frequent harmonic tetragrams (four consecutive chords) per collection (e.g.\ \texttt{IV$\cdot$I$\cdot$V$\cdot$vi}). The gradient of black and white shows stronger prevalence in darker cells, and integer annotations report absolute counts. The left panel is sorted using Suno as the reference; the right panel uses Udio as the reference.}
    \label{fig:bpe_trigrams_suno}
\end{figure*}

We also computed the ``Unique Tokens Ratio,'' which reflects how many different segment types appear in a song. Lower values suggest repetitive forms; higher values indicate more varied structures. This metric was lowest in Suno ($0.44\pm0.18$) and highest in Udio ($0.58\pm0.20$), with Lastfm ($0.51\pm0.19$) in between.

% \subsubsection{Summary}

% The results from our BPE analysis suggest several insights relevant to broader questions about the formal structures and creative capacities of AI-generated music.

% First, we observe that Suno exhibits a strong preference for simple, repetitive structures, particularly repeated verse-chorus cycles. This convergence toward a reduced formal vocabulary likely reflects biases in generation objectives optimized for familiarity and accessibility, as well as the short song snippets to be extended later via the web interface. Udio shows similar structural tendencies but with merge frequencies closer to Lastfm. A key confound for both platforms is song duration: Suno tracks average 175 seconds (approximately 2.9 minutes), Udio tracks average 143 seconds (approximately 2.4 minutes), and Lastfm tracks average 261 seconds (approximately 4.4 minutes), making AI-generated tracks roughly half to two-thirds the length of their human-composed counterparts. The reduced structural diversity observed in Suno and Udio may therefore reflect not only a bias toward formulaic composition but also a constraint imposed by shorter generation lengths, as shown in Figure~\ref{fig:durations}.


\subsection{Axis of Awesome Tendencies}
\label{sec:axis_tendencies}


\rev{We next measure how strongly each collection leans on the Axis of Awesome, counted two ways: the four chords in any order, and the literal loop (the chords in their cyclic order). Table~\ref{tab:axis_classification} reports, for each measure, the proportion of songs whose axis rate, the share of their four-chord windows that match (Section~\ref{sec:axis}), exceeds a range of thresholds.}

\begin{table}[ht]
\centering
\footnotesize
\caption{\rev{Proportion of songs (\%) per collection whose axis rate exceeds each threshold, under the two measures: the four chords in any order (72 patterns) and the loop only (12 rotations). Collections: Lastfm ($n=18{,}740$), Suno ($n=18{,}824$), Udio ($n=18{,}691$), counting songs with at least one four-chord window.}}
\label{tab:axis_classification}
\begin{tabular}{lrrrrrr}
\toprule
\textbf{Thresh.}
  & \multicolumn{3}{c}{\textbf{\rev{Any order (72)}}}
  & \multicolumn{3}{c}{\textbf{\rev{Loop only (12)}}} \\
\cmidrule(lr){2-4}\cmidrule(lr){5-7}
 & \rev{L} & \rev{S} & \rev{U} & \rev{L} & \rev{S} & \rev{U} \\
\midrule
\rev{$>$0\%}     & \rev{26.9} & \rev{47.1} & \rev{27.1} & \rev{10.5} & \rev{27.5} & \rev{10.6} \\
\rev{$\geq$1\%}  & \rev{24.9} & \rev{46.1} & \rev{26.2} & \rev{9.3}  & \rev{26.4} & \rev{9.8} \\
\rev{$\geq$5\%}  & \rev{13.4} & \rev{35.1} & \rev{16.0} & \rev{4.0}  & \rev{17.5} & \rev{5.0} \\
\rev{$\geq$10\%} & \rev{8.6}  & \rev{27.9} & \rev{10.5} & \rev{2.8}  & \rev{13.6} & \rev{3.3} \\
\rev{$\geq$25\%} & \rev{3.5}  & \rev{16.1} & \rev{4.4}  & \rev{1.3}  & \rev{8.2}  & \rev{1.6} \\
\rev{$\geq$50\%} & \rev{0.8}  & \rev{6.5}  & \rev{1.3}  & \rev{0.3}  & \rev{3.8}  & \rev{0.6} \\
\bottomrule
\end{tabular}
\end{table}

\rev{At the $\geq$10\% threshold, counting the four chords in any order, $8.6\%$ of Lastfm, $10.5\%$ of Udio, and $27.9\%$ of Suno songs qualify as axis-based; counting only the loop, the rates are $2.8\%$, $3.3\%$, and $13.6\%$ (Table~\ref{tab:axis_classification}). Figure~\ref{fig:axis_distribution} shows the distribution of per-song rates behind these counts. Given the collection sizes, the $95\%$ confidence interval of any proportion in Table~\ref{tab:axis_classification} is below $\pm0.8$ percentage points, so the differences between Suno and the other collections far exceed sampling uncertainty. Both measures confirm the findings of \citet{richards2017tonal} on the cultural ubiquity of this progression and extend them to show that Suno reproduces it at a substantially higher rate than both the human reference corpus and Udio, with the lead widening when only the loop counts (a Suno-to-Lastfm ratio of $3.2\times$ rising to $4.8\times$).}

\begin{figure}[!ht]
    \centering
    \includegraphics[width=\linewidth]{figures/tetragram_axis_distribution.png}
    \caption{\rev{Per-song axis rate under the two measures, for songs with at least 10\% axis content. (a)~Any order: a four-chord window counts when its chords form an axis set in any order (72 patterns across the three sets). (b)~The loop: a window counts only when it is one of the 12 rotations of the three loops (four per set). Curves are kernel-density estimates; each legend percentage is the share of that collection's songs above the 10\% threshold under that measure.}}
    \label{fig:axis_distribution}
\end{figure}

\rev{The rates above count how many songs lean on the axis; they do not say, when a collection reaches for these four chords, whether it plays the loop or rearranges them. Figure~\ref{fig:axis_rotation} answers that, splitting the four-chord windows built on the axis chords into the loop itself and the same chords in a different order. Panel~(a) gives each as a share of all four-chord windows: Suno places 12.3\% on the axis chords (3.4\% Lastfm, 4.5\% Udio), most of it in the loop. Panel~(b) is the loop's share of that usage: 54.3\% for Suno, against 36.3\% for Lastfm and 38.5\% for Udio. When Suno uses these chords, it usually plays the loop rather than a rearrangement.}

\begin{figure}[!ht]
    \centering
    \includegraphics[width=\linewidth]{figures/axis_rotation_breakdown.png}
    \caption{\rev{Loop versus reordering. Among four-chord windows built on the axis chords, the loop (the chords in cyclic order) is separated from the same chords in a different order. (a)~Each as a share of all four-chord windows, loop (solid) and reordering (hatched); the number inside the solid segment is the loop share alone, and the number above each bar is the combined total. (b)~The loop's share of those axis-chord windows. The dashed line marks 50\%.}}
    \label{fig:axis_rotation}
\end{figure}


% ------------------------------------------------------------------
\section{Discussion}
\label{sec:discussion}

% Explain this is an attempt at analyzing the musical aspects and musical form of AI vs Humans' music by implementing a state-of-the-art MIR tool. 
% Discuss the question: 

We now discuss the three research questions about AI-generated music at the corpus scale and then draw out two cross-cutting implications: outcome homogenization and creative agency. % (Section~\ref{sec:homogenization}).

\rev{We took several steps to characterize potential selection bias in the corpus. A popularity-filtered dataset would yield a distinctive median difference relative to an unfiltered baseline; the equivalence observed in the engagement distributions (Section~\ref{sec:corpora}) is instead consistent with unfiltered acquisition from each platform's front page, and no ranking, curation, or engagement threshold was applied to either collection. Although the Udio query nominally ranked tracks by 24-hour play count, those plays operate at very low absolute magnitudes (median: $1$ like, $4$ plays), so the ranking behaves as a near-uniform sample over recently uploaded content rather than as a popularity filter. Nevertheless, the} internal curation logic of Suno's \textit{New Songs} playlist and Udio's search index is undisclosed. Differences reported here should therefore be read as differences between the music surfaced by each platform, not strictly between unconstrained generative-model outputs.

% \begin{description}
%   \item[RQ1] How reliably can current MIR tools analyze harmonic structure, vocabulary, and musical form in AI-generated music?
%   \item[RQ2] To what extent does the harmony of AI-generated music relate to human-composed music?
%   \item[RQ3] What structural similarities or differences emerge when now comparing the and human-composed music?
% \end{description}

\subsection{RQ1: Applying MIR Tools to Large Music Corpora}
\label{sec:disc_mir}

The chord estimator and structural model adopted in this study 
% (Sections~\ref{sec:ace_pipeline} and~\ref{sec:form_pipeline}) 
address the principal methodological concerns for corpus-scale analysis, but challenges remain that should be acknowledged. For chord estimation, even state-of-the-art models exhibit lower accuracy on complex chord types (sevenths, extensions, suspensions, inversions) than on triads~\citep{pauwels2013evaluating}, and reported performance is sensitive to evaluation vocabulary and chord mapping strategies. \rev{The conversion from chord labels to Roman numerals depends on a single global key estimated per track, which can be wrong under modulation, strong modal mixture, or relative-major/minor ambiguity, and then propagates to every Roman numeral in the track. The same estimator runs identically on all three collections, so any such error is shared and cannot create the between-collection differences we report.} Furthermore, consonance-ACE assumes enharmonic equivalence: sharp ninths, sharp elevenths, and flat thirteenths are encoded as their enharmonic equivalents (flat third, flat fifth, and flat sixth, respectively), and therefore do not appear as distinct labels in the chord vocabulary.

This study operates at the level of chord-label distributions and n-gram sequential patterns, which is the appropriate granularity for corpus-scale comparison across $60,014$ tracks. Voice leading, harmonic rhythm, and chord-voicing detail are not captured in any chord-symbol-based corpus study and fall outside the present scope.

For structural analysis, the All-In-One model provides semantic section labels (verse, chorus, bridge, intro, outro) unavailable from unsupervised methods. Its reported scores on the Harmonix Set ($HR.5F = 0.660$ for segmentation, $PWF = 0.738$ for structure labeling) remain below estimated human inter-rater agreement~\citep{nieto2020audio}, so a degree of error is present in the annotations. However, since the same model is applied uniformly across all three collections, these errors affect each corpus equally. Our analysis targets distributional tendencies across $60,014$ tracks; systematic differences between collections therefore reflect genuine structural tendencies rather than tool-specific artifacts.

\rev{Platform genre tags record the surface aesthetic a user asks for (timbre, mood, instrumentation), not harmonic function, and the two are decoupled: the `Jazz' tag on Suno and Udio returns mostly ambient instrumental tracks, without the ii-V-I cadential vocabulary that defines the style. The per-genre panels (Figure~\ref{fig:trigrams_per_style}) therefore report what each platform files under a label, not a fixed musicological category, and our corpus-level contrasts do not depend on this: they are computed across all tracks regardless of tag.}

Despite these caveats, the combination of a chord estimator with a semantically aware structure analyzer represents a meaningful advance over the triad-level, label-free pipelines that were the practical default for large-scale audio analysis until recently.

\subsection{RQ2: Harmonic Relationship Between AI-Generated and Human-Composed Music}
\label{sec:disc_harmony}

The four metrics used in this section capture different dimensions of harmonic practice and should be read accordingly. Chord-quality distributions (Figure~\ref{fig:vocabulary_distribution}) describe the harmonic palette, which chord types appear and how often. Trigram frequencies and cross-collection correlations describe local sequential behavior, in which three-chord successions recur and at what relative rates. \rev{Axis rates measure use of a canonical four-chord set, reported both as family membership (any ordering) and as the literal loop (the canonical rotation order).} Tetragram frequencies describe recurring four-chord patterns and how often they recur within a collection. These dimensions are complementary but independent: a high trigram-frequency correlation between two collections, such as the $r = 0.952$ shared by Udio and Lastfm, indicates distributional similarity in local harmonic motion but says nothing about whether those progressions occupy the same formal positions within a song or occur at comparable rates of harmonic change.

Three strands of evidence converge on a consistent picture of how AI-generated harmony relates to human-composed music.

At the level of chord-quality vocabulary, the three collections are broadly similar: major triads dominate, followed by minor triads, power chords, seventh chords, dominant sevenths, and suspended chords (Figure~\ref{fig:vocabulary_distribution}). This indicates that current generative platforms reproduce the surface-level harmonic palette of popular music with fidelity. The differences emerge not in what chord types appear, but in how they are sequenced.

Trigram frequency analysis reveals that Suno concentrates its harmonic motion on a narrow subset of Western pop progressions. Across all 53~canonical progressions that pass the enrichment filters (FDR-adjusted $p \ll 0.001$, RR~$\geq 1.3$), Suno is overrepresented at frequency ratios ranging from $1.3\times$ to $3.3\times$ relative to Lastfm. The highest-frequency 15 of these 53~progressions, listed in Panel~A of Table~\ref{tab:suno_enrichment}, concentrate the over-representation at $2.1\times$--$3.3\times$, while the remaining lower-frequency canonical progressions contribute ratios in the $1.3\times$--$2.1\times$ range. The cumulative effect across all 53 enriched progressions is substantial: Suno's harmonic vocabulary is contracted to $6,770$ unique trigram types, compared to $8,056$ in Udio and $9,838$ in Lastfm.

\rev{At the four-chord level, the Axis of Awesome analysis (Section~\ref{sec:axis_tendencies}) sharpens the same contrast. Suno is axis-based far more often than either human-composed music or Udio, and the gap widens when matching is restricted to the literal loop rather than any ordering of the four chords: the order-sensitive test does not soften Suno's lead but increases it at every threshold. Suno does not merely share the axis vocabulary; it plays the loop in its canonical order far more often than human music does.}

Udio's harmonic distributions, by contrast, are statistically closer to those of Lastfm across all three measures. Cross-collection trigram frequency correlations confirm this: Udio and Lastfm share $r = 0.952$, while Suno correlates with Lastfm at only $r = 0.441$. The harmonic tetragram frequencies (Figure~\ref{fig:bpe_trigrams_suno}) reinforce this at the four-chord level: Udio and Lastfm share a similar tetragram distribution ($r = 0.968$) while Suno diverges toward a smaller set of high-frequency patterns ($r = 0.779$ with Lastfm). Udio features minor-key functional progressions (e.g., \texttt{bVI-bVII-i}, \texttt{i-bVI-III}, \texttt{i-bVII-bVI}) among its most frequent trigrams (see Figure~\ref{fig:top_30_trigrams}). Progressions of this type, particularly the \texttt{bVI-bVII-i} pattern, have been identified as characteristic harmonic idioms of rock music, functioning as Aeolian cadences with subdominant--dominant--tonic implications distinct from conventional major-minor tonality~\citep{biamonte2010triadic}. Lastfm's top trigrams are dominated by functional major progressions (\texttt{IV-I-V}, \texttt{IV-V-I}, \texttt{V-IV-I}), with minor-mode patterns present but appearing further down the frequency ranking (see Figure~\ref{fig:top_30_trigrams}).

\subsection{RQ3: Structural Similarities and Differences}
\label{sec:disc_structure}

The n-gram analysis of section-label sequences and the Unique Tokens Ratio provide complementary views of formal diversity across the three collections.

Suno exhibits the strongest convergence toward simple, repetitive structures. Verse-chorus cycles (ABAB patterns) dominate its most frequent section-label tetragrams, and its Unique Tokens Ratio is the lowest of the three collections ($0.44 \pm 0.18$), indicating that the average Suno track draws from fewer distinct section types than either human-composed or Udio-generated music. Udio, by contrast, shows the highest formal diversity ($0.58 \pm 0.20$), with Lastfm occupying an intermediate position ($0.51 \pm 0.19$). The section-label tetragram heatmap (Figure~\ref{fig:musical_form_heatmap}) confirms that Udio's tetragram frequencies are closer to Lastfm than to Suno. Of particular note, the second most frequent pattern in both AI collections, \texttt{verse·chorus·inst~$\rightarrow$~verse} (Suno: $2,103$; Udio: $1,599$), is entirely absent from Lastfm. Both systems frequently produce instrumental sections as structural connectors within verse-chorus cycles, a formal strategy with no parallel in the human corpus. 

The normalized musical form trigram profiles per genre (Figure~\ref{fig:musical_form}) reveal a consistent pattern across styles. Both Suno and Udio show a reduced range of formal options relative to Lastfm, with intro-verse-chorus sequences markedly under-represented and solo sections largely absent. This is consistent with how both platforms operate: users generate a song segment that begins directly with a verse-chorus cycle and can later be extended via the web interface. The result is that intro sections, which precede the core structure, and solo sections, which typically appear later in human-composed music, are rarely generated in the initial output. The same genre-ontology caveat raised for Figure~\ref{fig:trigrams_per_style} applies here: the Jazz polygon in Figure~\ref{fig:musical_form} should not be read as a direct cross-collection comparison of jazz formal conventions, since the tracks grouped under that label differ substantially in character across the three corpora.

A key confound is song duration. 
Suno  and Udio tracks average $175$ and $143$ seconds respectively, against $261$ seconds for Lastfm (Figure~\ref{fig:durations}). 
Shorter tracks have fewer opportunities for structural elaboration (bridges, instrumental breaks, key changes), so the reduced diversity observed in Suno and Udio may partly reflect a generation-length constraint rather than a purely compositional bias. Udio's duration distribution is notably bimodal: a sharp spike below one minute reflects tracks published at the platform's default clip length of 32 seconds without extension, while the secondary peak near $130$ seconds corresponds to tracks generated at the platform's longest available clip length.

\begin{figure}[!htb]
    \centering
    \includegraphics[width=\linewidth]{figures/durations.png}
    \caption{Distribution of song durations across the three collections (Lastfm, Suno, Udio). The x-axis represents duration in seconds, and the y-axis indicates the number of songs per 10-second bin.}
    \label{fig:durations}
\end{figure}

A further caveat concerns formal labels and musical function. The All-In-One model assigns labels based on acoustic properties, but a verse-chorus alternation in a generated track may represent static loop repetition rather than the narrative-arc form those labels imply in conventional songwriting~\citep{nieto2020audio}. Distinguishing conventional pop form from non-teleological loop structures would require additional features such as melodic contour, rhythmic analysis, or lyric structure.

% 5) Why is it not possible to prompt harmonies into Suno or Udio?
\subsection{Outcome Homogenization and Creative Agency}
\label{sec:homogenization}
The harmonic compression documented in RQ2 and the structural narrowing documented in RQ3 both point toward a form of outcome homogenization, consistent with the concept of algorithmic monoculture, where shared components across AI systems increase the likelihood of convergent outputs~\citep{bommasani2022picking}. Suno exemplifies this pattern, converging on a statistical peak of canonical pop progressions (Figure~\ref{fig:pop_prevalence}, Table~\ref{tab:suno_enrichment}), also visible in the most common harmonic tetragrams (Figure~\ref{fig:bpe_trigrams_suno}). The output does not reflect the variety of harmonic practice or formal organization found in human-composed music.
Several mechanisms may contribute to this convergence. First, training data bias amplification, in which generative models learn to reproduce only the most frequent patterns in their training corpora. Since contemporary pop music dominates streaming platforms and likely constitutes a significant portion of Suno's training data, these harmonic patterns are overrepresented in its outputs. Second, optimization toward perceptual plausibility encourages models to generate music that sounds ``familiar'' to broad audiences, reinforcing existing harmonic conventions rather than exploring innovative progressions. Third, the lack of harmonic-prompting capabilities in current platforms means that users cannot directly specify harmonic content, instead relying on the models' learned biases toward common progressions. Fourth, Suno's developers may intentionally be targeting users who prefer a narrower spectrum of music, focusing on commercial modern pop styles.
Our data cannot determine which of these factors is primary, but the contrast with Udio is informative. Udio produces more diverse harmonic and structural outputs than would be expected from a presumably comparable training distribution, suggesting that model architecture and generation strategy play a role beyond the training data alone.
These platforms promise to make music creation widely accessible, but generating convincing audio recordings does not necessarily entail preserving musical knowledge or variety. A restricted harmonic vocabulary is not in itself a compositional deficiency: twelve-bar blues, punk, and minimal techno are all built on narrow chord sets by deliberate stylistic choice. The data show that Suno applies a single harmonic template across all stated genres, a platform-level convergence that is independent of user intent. This is consistent with Suno's own positioning: the platform's leadership has described it explicitly as gamifying musical creation.\endnote{Adam Neely: Suno, `AI Music, and the Bad Future'. \url{https://www.youtube.com/watch?v=U8dcFhF0Dlk}. Last accessed May 05, 2026} This raises questions about whether these tools genuinely support creative agency among musicians with any degree of musical knowledge~\citep{wiggins2007compositional, wiggins2016musical, lin2023promptsexploringdesignspace, tan2024we}. The prevalence of simple, repetitive formal structures (as revealed by our section-label analysis shown in Figure~\ref{fig:musical_form_heatmap}) further suggests that current platforms may be optimized to generate familiar, commercially viable content rather than to support experimentation or skill acquisition. This optimization strategy may serve commercial interests, but it limits the creative potential of AI as a compositional tool.




% ------------------------------------------------------------------
\section{Conclusion}
\label{sec:conclusion}

Across all three research questions, a consistent picture emerges. Applied uniformly to $60,014$ tracks, current MIR tools reveal that the two AI platforms under study diverge from each other as substantially as they diverge from human-composed music. Suno concentrates its harmonic output on a narrow subset of Western pop progressions, overrepresented at \rev{$1.3$ to $3.3$} times the rate found in the human reference corpus, and its formal repertoire is dominated by simple verse-chorus cycles with little structural variety. Udio, by contrast, exhibits harmonic and structural distributions that are statistically closer to those of human-composed music across every measure examined: trigram frequency correlations, axis-family rates, tetragram frequency profiles, and formal diversity all converge on the same ordering, Udio closer to Lastfm, Suno further away. Because both platforms draw on popular music at a corpus-level scale, this divergence implicates platform-level choices, model architecture, generation strategy, and optimization objectives as primary determinants of the harmonic and formal character of AI-generated music. This narrowness emerges despite diverse user intent. The analysis from \rev{\citet{casini2025data}} shows that Suno and Udio users actively prompt with distinct style descriptions, alongside explicit decade- and era-conditioning, indicating that the contracted harmonic vocabulary documented here is a property of the model's generation rather than of homogeneous user demand or restricted training exposure.

The most pressing methodological extension is to improve functional harmonic labeling to handle modulation, modal mixture, and key ambiguity, thereby sharpening analyses currently dependent on a single global key estimate per track. Tracking how these patterns evolve across successive model versions would further test whether the divergence documented here is a stable property of each platform or a moving target.

More broadly, this work demonstrates that fully automated, open-source MIR pipelines can now support corpus-scale harmonic and formal analysis at a scale and vocabulary richness that were not previously practical. The combination of rich-vocabulary chord estimation and semantic section labeling, applied here to $60,014$ tracks, offers a replicable methodology for computational musicology and MIR as AI-generated music continues to grow in both volume and diversity.

% ------------------------------------------------------------------
\section{Reproducibility}
\label{sec:reproducibility}
We release the complete dataset (60,014 JSON files with beat-level chord annotations, segmentation boundaries, style metadata, and audio links) and all analysis code under an MIT license.\endnote{\url{https://anonymous.4open.science/r/Aimir_Structure_Analysis/README.md}}.

% ------------------------------------------------------------------
\section{Ethics and consent}

Although the full ramifications of AI-generated material remain uncertain, the sheer speed and volume at which these systems can create new works demand scholarly attention \citep{Sturm2024a,pelly2025mood}. Because our study centers on music, we gathered examples from the commercial services {\em Suno} and {\em Udio}. Doing so raises a set of ethical issues. Even when a site's terms of use do not appear to permit automated collection, the legitimacy of repurposing that material for research, particularly under a "fair use" rationale, deserves scrutiny. As a publicly funded team investigating AI music and its broader consequences, we rely on the text-and-data-mining exception laid out in the EU Copyright Directive 2019/790 \citep{EUDirective2019} to justify our data acquisition.
%%%%%%%%%%%%%%%%%%%%%%%%%%%%%%%%%%%%%%%%%%%%%%%%%%%%%%%%%%%%%%%%%%%%%%%%%%%%%%%%
% Please do not touch.
% Print Endnotes
\IfFileExists{\jobname.ent}{
   \theendnotes
}{
   %no endnotes
}
%%%%%%%%%%%%%%%%%%%%%%%%%%%%%%%%%%%%%%%%%%%%%%%%%%%%%%%%%%%%%%%%%%%%%%%%%%%%%%%%

\section*{Acknowledgements}

% This paper is an outcome of a project that has received funding from the European Research Council (ERC) under the European Union's Horizon 2020 research and innovation programme (MUSAiC, Grant Agreement No. 864189).

\section*{Contributions}
% XXX - DD
% YYY - LCV
% ZZZ - BS
% WWW - 
XXX did ... YYY did ... 
%%%%%%%%%%%%%%%%%%%%%%%%%%%%%%%%%%%%%%%%%%%%%%%%%%%%%%%%%%%%%%%%%%%%%%%%%%%%%%%%
% Bibliography
%%%%%%%%%%%%%%%%%%%%%%%%%%%%%%%%%%%%%%%%%%%%%%%%%%%%%%%%%%%%%%%%%%%%%%%%%%%%%%%%

% For bibtex users:
\bibliography{TISMIRtemplate}

% For non bibtex users:
%\begin{thebibliography}{citations}
%
%\bibitem {Author:00}
%E. Author.
%``The Title of the Conference Paper,''
%{\it Proceedings of the International Symposium
%on Music Information Retrieval}, pp.~000--111, 2000.
%
%\bibitem{Someone:10}
%A. Someone, B. Someone, and C. Someone.
%``The Title of the Journal Paper,''
%{\it Journal of New Music Research},
%Vol.~A, No.~B, pp.~111--222, 2010.
%
%\bibitem{Someone:04} X. Someone and Y. Someone. {\it Title of the Book},
%    Editorial Acme, Porto, 2012.
%
%\end{thebibliography}

\end{document}